% Generated by scripts/tikz_reviews.py from the figure manifests.
\ReviewFigure{Convolution with a box window}
{figures/2-fourier/convolution.png}
{figures/tikz/fourier-convolution.pdf}
{Convolution; equation \eqref{eq-convol-1d}.}
{The shaded interval is [x-epsilon,x+epsilon]. The original box is unnormalized, so its convolution is an integral over that interval, not an average. The drawn signal is illustrative; no numerical data are asserted.}
{fourier-convolution}
\ReviewFigure{Convolution becomes multiplication}
{figures/2-fourier/convolution-fourier.png}
{figures/tikz/fourier-convolution-fourier.pdf}
{Convolution theorem; equation \eqref{eq-convol-fourier}.}
{The diagram uses the chapter convention F(f*g)=f-hat times g-hat. Both vertical arrows apply the Fourier transform, making the square commute; the original reverse Fourier arrow is equivalently written in the forward direction.}
{fourier-convolution-fourier}
\ReviewFigure{Cardinal splines by convolution}
{figures/2-fourier/splines.pdf}
{figures/tikz/fourier-splines.pdf}
{Convolution; Figure \ref{fig-splines}.}
{Interpretation to check: The box, triangle, and quadratic spline are exact centered cardinal B-splines B0, B1=B0*B0, and B2=B1*B0. The original final handwritten symbol is unclear; explicit degree and convolution labels replace it.}
{fourier-splines}
\ReviewFigure{Translation equivariance}
{figures/2-fourier/translation-inv.png}
{figures/tikz/fourier-translation-inv.pdf}
{Translation-invariant operators; T\_tau f(x)=f(x-tau).}
{The original square is made explicit with T\_tau f=f(.-tau). The bottom-right value is H(T\_tau f)=T\_tau(Hf); this is equivariance, called translation invariance in the chapter.}
{fourier-translation-inv}
\ReviewFigure{Fourier modes align at the Dirac comb}
{figures/2-fourier/poisson-distrib.png}
{figures/tikz/fourier-poisson-distrib.pdf}
{Poisson formula and distributions; finite Dirichlet kernel.}
{Three real Fourier modes cos(n omega), n=1,2,3, are shown. They all equal one at multiples of 2pi. The distributional comb identity is stated separately; a finite sum is not drawn as literal Dirac masses.}
{fourier-poisson-distrib}
\ReviewFigure{One radix-two FFT step}
{figures/2-fourier/fft.pdf}
{figures/tikz/fourier-fft.pdf}
{Discrete Fourier transform; recursion immediately before Figure \ref{fig-fft}.}
{The original butterfly wiring is replaced by an equivalent block diagram for the chapter's decimation-in-frequency recursion. Sum and difference are formed from the two halves, twiddle factors multiply the difference before its half-size FFT, and outputs are interleaved into even and odd frequencies.}
{fourier-fft}
\ReviewFigure{Spatial zero padding refines frequency sampling}
{figures/2-fourier/padding-spacial.png}
{figures/tikz/fourier-padding-spatial.pdf}
{Fourier approximation via spatial zero padding; Figure \ref{fig-padding-spacial}.}
{Interpretation to check: The spatial extent is T=Q/N while the physical sample step remains 1/N. The frequency interval is bounded by the unchanged Nyquist frequency pi N, and its step is 2pi/T. This corrects the handwritten reciprocal Nyquist labels. Curves are schematic and do not assert empirical sample values. The spectrum is shown by its magnitude; the two continuous curve profiles are schematic rather than an asserted exact transform pair.}
{fourier-padding-spatial}
\ReviewFigure{Spectral zero padding interpolates the samples}
{figures/2-fourier/padding-fourier.png}
{figures/tikz/fourier-padding-fourier.pdf}
{Fourier interpolation via spectral zero padding; Figure \ref{fig-padding-fourier}.}
{Interpretation to check: The retained signed-frequency coefficients are padded with zeros, and the inverse transform is multiplied by Q/N. The illustrative trigonometric polynomial and coefficient magnitudes are consistent for N=5 and Q=15; dark markers identify the five original samples and smaller colored markers the finer grid. These are analytic example values, not measurements. Odd N avoids the even-length Nyquist splitting convention.}
{fourier-padding-fourier}
\ReviewFigure{A plane wave and its normal direction}
{figures/2-fourier/wave.png}
{figures/tikz/fourier-wave.pdf}
{Fourier in multiple dimensions; tensor products of complex exponentials.}
{The parallel lines are constant-phase sets, normal to the wave vector omega. They indicate geometry rather than a propagation direction. Integer wave vectors apply on the periodic domain; the general figure uses omega in R squared.}
{fourier-wave}
\ReviewFigure{A square with opposite edges identified}
{figures/2-fourier/torus.png}
{figures/tikz/fourier-torus.pdf}
{Multidimensional torus; (R / 2pi Z) squared.}
{Opposite sides of the square carry matching oriented arrows. Identifying each pair produces the two-dimensional torus. The pictured doughnut is a topological visualization, not an isometric embedding of the flat quotient.}
{fourier-torus}
\ReviewFigure{A periodic two-dimensional grid}
{figures/2-fourier/torus-discr.png}
{figures/tikz/fourier-torus-discrete.pdf}
{Discrete multidimensional Fourier basis; product of cyclic index sets.}
{A four-by-four representative grid illustrates the two periodic coordinate directions. A step beyond the last displayed index wraps back to zero; the first and last displayed vertices remain distinct. The identifications occur after N1 or N2 steps, so the rectangle has no physical boundary. The grid size is illustrative.}
{fourier-torus-discrete}
\ReviewFigure{Heat diffusion is Gaussian convolution}
{figures/2-fourier/heat.png}
{figures/tikz/fourier-heat.pdf}
{Fourier multipliers and differential equations; Figure \ref{fig-heat}.}
{The left curves show an illustrative signal with two Fourier modes at time zero and after positive diffusion times. The right curves are Gaussian heat kernels for the same two times: increasing time broadens the kernel and lowers its peak. The normalization and variance follow the chapter heat equation.}
{fourier-heat}
\ReviewFigure{Continuum and discrete Laplacian spectra}
{figures/2-fourier/finite-diff.png}
{figures/tikz/fourier-finite-difference-spectrum.pdf}
{Finite domain and discretization; equations \eqref{eq-disc-diff-2} and \eqref{eq-fft-lapl}.}
{The original positive curves are explicitly labeled as negative eigenvalues divided by N squared. The continuum curve is (2pi k/N) squared, reaching pi squared at Nyquist; the discrete curve is 4 sin squared(pi k/N), reaching four. Actual Laplacian eigenvalues are nonpositive.}
{fourier-finite-difference-spectrum}
\ReviewFigure{A small geodesic ball on a surface}
{figures/2-fourier/laplacian-surf.png}
{figures/tikz/fourier-laplacian-surface.pdf}
{Laplace--Beltrami operator and the small-ball mean-value expansion.}
{The sketch is a local surface patch containing a geodesic ball centered at x. The ellipse represents the projected ball, not a claim that geodesic balls are Euclidean circles. The displayed expansion retains the necessary epsilon-squared scaling and dimension-dependent constant.}
{fourier-laplacian-surface}
\ReviewFigure{Colatitude and azimuth on the sphere}
{figures/2-fourier/spherical.png}
{figures/tikz/fourier-spherical-coordinates.pdf}
{Spherical harmonics; x=(sin theta cos phi,sin theta sin phi,cos theta).}
{Interpretation to check: Theta is measured from the positive north axis, matching the chapter colatitude convention. The original angle appears measured from the equatorial projection, which would instead be latitude; this ambiguity is deliberately resolved from the formula. Phi is the equatorial azimuth.}
{fourier-spherical-coordinates}
\ReviewFigure{A weighted graph and its local Laplacian}
{figures/2-fourier/graph.png}
{figures/tikz/fourier-weighted-graph.pdf}
{Graph Laplacian; Delta=W-D.}
{Interpretation to check: The graph patch preserves the intended local adjacency picture, highlighting an edge between vertices i and j with weight w\_ij. Unreadable numerical marks are not assigned values. The Laplacian sign matches the chapter nonpositive convention.}
{fourier-weighted-graph}
